\documentclass{article}
\usepackage{amsmath}
\usepackage{amssymb}
\usepackage{booktabs}
\usepackage{multirow}
\usepackage[utf8]{inputenc}
\usepackage{graphicx}
\usepackage{caption}
\usepackage[a4paper, margin=0.8in]{geometry}
\everymath{\displaystyle}
\usepackage{xcolor}
\usepackage{listings}
\definecolor{grey}{rgb}{0.8,0.8,0.8}
\definecolor{darkgreen}{rgb}{0,0.3,0}
\definecolor{darkblue}{rgb}{0,0,0.3}
\def\lstbasicfont{\fontfamily{pcr}\selectfont\footnotesize}
\renewcommand{\lstlistingname}{}
\lstset{%
    numbers=none,
    numberstyle=\small,%
    showstringspaces=false,
    showspaces=false,%
    tabsize=4,%
    frame=lines,%
    basicstyle={\footnotesize\lstbasicfont},%
    keywordstyle=\color{darkblue}\bfseries,%
    identifierstyle=,%
    commentstyle=\color{darkgreen},%\itshape,%
    stringstyle=\color{black},%
    breaklines=true,%
    postbreak=\mbox{\textcolor{red}{$\hookrightarrow$}\space}
}
\lstloadlanguages{C,C++,Java,Matlab,Python,Mathematica}


\title{DOBI: Mathematical Formulation and Parameter Analysis}
\author{Xiuyuan Lu}
\date{}

\begin{document}

\maketitle

\section{Introduction}
The Dombi-Bonferroni (DOBI) method is a novel multi-criteria decision making (MCDM) technique that combines Dombi operations with Bonferroni mean operators. This document presents its mathematical formulation and analyzes the effects of its key parameters.

\section{Mathematical Formulation}

\subsection{Data Normalization}
The method begins with normalized decision matrix $X^*$ where:

\begin{equation}
x_{ij}^* = 
\begin{cases}
\frac{x_{ij}}{\max(x_j)} & \text{(benefit criterion)} \\
\frac{\min(x_j)}{x_{ij}} + \frac{\max(x_j)}{\max(x_j)} + \frac{\min(x_j)}{\max(x_j)} & \text{(cost criterion)}
\end{cases}
\end{equation}

\subsection{Additive Normalization (f-values)}
\begin{equation}
f_{ij} = \frac{x_{ij}^*}{\sum_{j=1}^n x_{ij}^*}
\end{equation}

\subsection{Dombi-Bonferroni Functions}

\subsubsection{Weighted Averaging Function (Z1)}
\begin{equation}
Z1_i = \frac{\sum_{j=1}^n x_{ij}^*}{1 + \left[\frac{\sum_{j=1}^n \sum_{\substack{k=1 \\ k \neq j}}^n \frac{w_j w_k}{1-w_j} \left(\psi_1 \left(\frac{1-f_{ij}}{f_{ij}}\right)^\zeta + \psi_2 \left(\frac{1-f_{ik}}{f_{ik}}\right)^\zeta\right)}{\psi_1 + \psi_2}\right]^{1/\zeta}}
\end{equation}

\subsubsection{Weighted Geometric Function (Z2)}
\begin{equation}
Z2_i = \frac{\sum_{j=1}^n x_{ij}^*}{1 + \left[\frac{\sum_{j=1}^n \sum_{\substack{k=1 \\ k \neq j}}^n \frac{w_j w_k}{1-w_j} \left(\psi_1 \left(\frac{f_{ij}}{1-f_{ij}}\right)^\zeta + \psi_2 \left(\frac{f_{ik}}{1-f_{ik}}\right)^\zeta\right)}{\psi_1 + \psi_2}\right]^{1/\zeta}}
\end{equation}

\subsection{Integrated DOBI Value}
\begin{equation}
\Re_i = \frac{Z1_i + Z2_i}{1 + \left[0.5\left(\frac{1-Z1_i}{Z1_i}\right)^\delta + 0.5\left(\frac{1-Z2_i}{Z2_i}\right)^\delta\right]^{1/\delta}}
\end{equation}

\subsection{Ranking}
\begin{equation}
\text{Rank}_i = \text{rank}(\Re_i, \text{ascending=False})
\end{equation}

\section{Parameter Analysis}

\begin{table}[h]
\centering
\caption{DOBI Parameters and Their Effects}
\begin{tabular}{lll}
\toprule
Parameter & Role & Impact \\
\midrule
$\psi_1,\psi_2$ & Bonferroni parameters & Control inter-criteria compensation \\
$\zeta$ & Dombi operator parameter & Adjusts aggregation behavior \\
$\delta$ & Integration parameter & Balances Z1 and Z2 contributions \\
\bottomrule
\end{tabular}
\end{table}

\subsection{Effects of Parameter Settings}

\subsubsection{$\psi_1$ and $\psi_2$ (Bonferroni Parameters)}
\begin{itemize}
\item \textbf{Symmetry}: When $\psi_1 = \psi_2$, equal weight to all interactions
\item \textbf{Compensation}: Higher values increase compensation effects
\item \textbf{Robustness}: Asymmetric values ($\psi_1 \neq \psi_2$) increase sensitivity to specific criteria
\end{itemize}

\subsubsection{$\zeta$ (Dombi Operator Parameter)}
\begin{equation}
\lim_{\zeta \to 0} \rightarrow \text{maximum-like operation}, \quad \lim_{\zeta \to \infty} \rightarrow \text{minimum-like operation}
\end{equation}

\begin{itemize}
\item Small $\zeta$: More sensitive to high values
\item Large $\zeta$: More sensitive to low values
\item $\zeta=1$: Linear compensation
\end{itemize}

\subsubsection{$\delta$ (Integration Parameter)}
\begin{itemize}
\item $\delta=1$: Equal treatment of Z1 and Z2
\item $\delta>1$: Emphasizes larger values
\item $\delta<1$: Reduces impact of extreme values
\end{itemize}

\section{Recommended Parameter Settings}

\begin{table}[h]
\centering
\caption{Parameter Setting Scenarios}
\begin{tabular}{lllll}
\toprule
Scenario & $\psi_1,\psi_2$ & $\zeta$ & $\delta$ & Application \\
\midrule
Default & 1,1 & 1 & 1 & General purpose \\
Robust & 1,1 & 2 & 0.5 & Noisy data \\
Dominance & 2,2 & 0.5 & 2 & Key criteria emphasis \\
\bottomrule
\end{tabular}
\end{table}

\section{Implementation Notes}
The Python implementation features:
\begin{itemize}
\item Direct matrix operations for efficiency
\item Flexible parameter configuration
\item Batch processing for temporal data
\end{itemize}

For sensitivity analysis, we recommend:

\begin{verbatim}
param_combinations = [
    (1, 1, 1, 1),  # Baseline
    (1, 2, 1, 1),  # Asymmetric psi
    (1, 1, 2, 1),  # Increased zeta
    (1, 1, 1, 2)   # Increased delta
]
\end{verbatim}

\section{Appendix}
Here is the source code for DOBI implementation in Python:
\lstinputlisting[language=Python, caption={}]{../../USdata/aggregating_DOBI_simplified.py}

\end{document}